%! no-theme
% three-architectures (ch07, rem-three-archs / Ex 7.2) — two binary switches (self-attention
% bidirectional vs causal) × (cross-attention absent vs present) give a larger four-cell switch space;
% three cells are the standard representation/generation patterns, while the fourth
% is a valid non-autoregressive fusion pattern. Drawing the 2×2 makes the scope distinction literal.
%   (bi, absent)   = Encoder-only   · BERT
%   (causal, absent)= Decoder-only  · GPT, LLaMA
%   (causal, present)= Encoder-decoder · T5, original transformer
%   (bi, present)  = two-stream bidirectional cross-attention fusion (valid, not autoregressive)
% Pedagogical mitigations: (i) each block carries an attention-mask icon (\tfmaskicon full/causal) so the
% row switch is SHOWN, not just named; (ii) the enc-dec cell shows BOTH a bidirectional-masked encoder and a
% causal-masked decoder (kills "enc-dec has no bidirectional part"); (iii) the 4th cell is valid but non-autoregressive.
% /authoring-figures standard: relative positioning + tokens, grayscale-safe (masks are shading; roles by
% label+border), line-free title/caption lanes.
\documentclass[tikz,border=10pt]{standalone}
\usepackage{transformer-figures}
\begin{document}
\begin{tikzpicture}
  \newcommand{\CW}{50mm}\newcommand{\CH}{40mm}\newcommand{\GAP}{7mm}

  % ══ the 2×2 cells (rows: bidirectional / causal · cols: cross-attn absent / present) ══
  \node[tfpanel,    minimum width=\CW, minimum height=\CH]                       (cA1) {}; % bi, absent
  \node[tfpanel,    minimum width=\CW, minimum height=\CH, right=\GAP of cA1]     (cB1) {}; % bi, present  → fusion
  \node[tfpanel,    minimum width=\CW, minimum height=\CH, below=\GAP of cA1]     (cA2) {}; % causal, absent
  \node[tfpanel,    minimum width=\CW, minimum height=\CH, below=\GAP of cB1]     (cB2) {}; % causal, present

  % ── cell A1 — Encoder-only ──
  \node[tfrole, font=\small\bfseries, anchor=north] at ([yshift=-2.5mm]cA1.north) {Encoder-only};
  \node[tfaxis, anchor=south] (a1o) at ([yshift=6mm]cA1.center) {pooled embedding};
  \node[tfstruct, minimum width=15mm, minimum height=12mm, align=center, font=\footnotesize] (a1b) at ([yshift=-1mm]cA1.center) {Enc\\$\times N$};
  \coordinate (a1m) at ([xshift=12mm]a1b.center); \tfmaskicon{full}{a1m}
  \node[tfaxis, anchor=north] (a1i) at ([yshift=-8.5mm]cA1.center) {tokens};
  \node[tfaxis, anchor=south] at ([yshift=1.5mm]cA1.south) {BERT};
  \draw[tfflow] (a1i) -- (a1b);   \draw[tfflow] (a1b) -- (a1o);

  % ── cell A2 — Decoder-only ──
  \node[tfrole, font=\small\bfseries, anchor=north] at ([yshift=-2.5mm]cA2.north) {Decoder-only};
  \node[tfaxis, anchor=south] (a2o) at ([yshift=6mm]cA2.center) {next token};
  \node[tfstruct, minimum width=15mm, minimum height=12mm, align=center, font=\footnotesize] (a2b) at ([yshift=-1mm]cA2.center) {Dec\\$\times N$};
  \coordinate (a2m) at ([xshift=12mm]a2b.center); \tfmaskicon{causal}{a2m}
  \node[tfaxis, anchor=north] (a2i) at ([yshift=-8.5mm]cA2.center) {prompt};
  \node[tfaxis, anchor=south] at ([yshift=1.5mm]cA2.south) {GPT, LLaMA};
  \draw[tfflow] (a2i) -- (a2b);   \draw[tfflow] (a2b) -- (a2o);

  % ── cell B2 — Encoder-decoder (BOTH masks; cross-attn couples them) ──
  \node[tfrole, font=\small\bfseries, anchor=north] at ([yshift=-2.5mm]cB2.north) {Encoder-decoder};
  \node[tfstruct, minimum width=12mm, minimum height=11mm, align=center, font=\scriptsize] (b2e) at ([xshift=-11.5mm,yshift=1mm]cB2.center) {Enc\\$\times N$};
  \node[tfstruct, minimum width=12mm, minimum height=11mm, align=center, font=\scriptsize] (b2d) at ([xshift=11.5mm,yshift=1mm]cB2.center) {Dec\\$\times N$};
  \coordinate (b2em) at ([yshift=-8mm]b2e.center); \tfmaskicon{full}{b2em}
  \coordinate (b2dm) at ([yshift=-8mm]b2d.center); \tfmaskicon{causal}{b2dm}
  \draw[tfflow] (b2e.east) -- (b2d.west) node[midway, tfoplabel, font=\tiny] {cross-\\attn};
  \node[tfaxis, anchor=south] at ([yshift=6mm]b2d.center) {next token};
  \node[tfaxis, anchor=south] at ([yshift=1.5mm]cB2.south) {T5, original transformer};

  % ── cell B1 — valid bidirectional fusion (not an autoregressive generator) ──
  \node[tfrole, font=\small\bfseries, anchor=north] at ([yshift=-2.5mm]cB1.north) {Bidirectional fusion};
  \node[tfstruct, minimum width=12mm, minimum height=11mm, align=center, font=\scriptsize] (b1e) at ([xshift=-11.5mm,yshift=1mm]cB1.center) {Source\\enc.};
  \node[tfstruct, minimum width=12mm, minimum height=11mm, align=center, font=\scriptsize] (b1f) at ([xshift=11.5mm,yshift=1mm]cB1.center) {Fusion\\stack};
  \coordinate (b1m) at ([yshift=-8mm]b1f.center); \tfmaskicon{full}{b1m}
  \draw[tfflow] (b1e.east) -- (b1f.west) node[midway, tfoplabel, font=\tiny] {cross-\\attn};
  \node[tfaxis, anchor=south] at ([yshift=6mm]b1f.center) {joint representation};
  \node[tfaxis, anchor=south] at ([yshift=1.5mm]cB1.south) {fusion, reranking};

  % ══ axis headers ══
  \coordinate (gtop) at ([yshift=3mm]cA1.north);
  \node[tfcap, anchor=south] at (cA1.north |- gtop) {absent};
  \node[tfcap, anchor=south] at (cB1.north |- gtop) {present};
  \coordinate (colmid) at ($(cA1.north)!0.5!(cB1.north)$);
  \node[tfrole, anchor=south] at ([yshift=7.5mm]colmid) {cross-attention};

  \coordinate (gleft) at ([xshift=-3mm]cA1.west);
  \node[tfcap, anchor=east] at (gleft |- cA1) {bidirectional};
  \node[tfcap, anchor=east] at (gleft |- cA2) {causal};
  \coordinate (rowmid) at ($(cA1.west)!0.5!(cA2.west)$);
  \node[tfrole, rotate=90, anchor=south] at ([xshift=-26mm]rowmid) {self-attention};

  % ══ title + mask legend + caption (line-free lanes) ══
  \node[tfrole, font=\normalsize\bfseries, anchor=south] at ([yshift=13mm]colmid) {Three standard patterns inside a larger switch space};

  \coordinate (botmid) at ($(cA2.south)!0.5!(cB2.south)$);
  \node[tfcap, align=center, anchor=north] at ([yshift=-5mm]botmid)
       {Row $=$ the self-attention of the output stack — encoder-decoder is causal because its decoder generates.\\
        Top-right is valid bidirectional fusion, but not an autoregressive generator.};
  \coordinate (legY) at ([yshift=-16mm]botmid);
  \coordinate (lm1) at ([xshift=-32mm]legY); \tfmaskicon{full}{lm1}
  \node[tfcap, anchor=west] at ([xshift=3.5mm]lm1) {attend all positions (bidirectional)};
  \coordinate (lm2) at ([xshift=24mm]legY); \tfmaskicon{causal}{lm2}
  \node[tfcap, anchor=west] at ([xshift=3.5mm]lm2) {attend earlier positions only (causal)};

\end{tikzpicture}
\end{document}
